\subsection{Higher even terminal lists with a nonunit ordinary anchor}\label{proof:R2-C036}

\paragraph{R36.1 Original object and exact ordinary theorem}

The original problem quantifies every finite tree \(U\) with at most five vertices of original degree \(\ge 3\). Every original edge direction must be fixed from \(U\) before every finite original probability \(\mu\) and real skew pointwise \(\lvert K\rvert \le 1\). The target is \(Z_D=E_\mu^V \prod(1+K)\le 1\). Every original arm, repeated arm, subdivision, all four independent positive core lengths, correlated joint law and noninjective tournament map remain in the problem.

\(At\) an original terminal \(t\) let the strictly ordered residual after genuine equal pairing be \(a_{1}<b_{1}<\ldots<a_{j}<b_{j},\allowbreak \allowbreak  a_{1}\ge 4,\allowbreak \allowbreak j\ge 1\). Let its actual \(F\) be \(\prod_i h_{a_{i}} f_{b_{i}}\). \(At\) receiver \(r\) there is an actual arm \(m\ge 3\) and an optional original root block \(P=\prod_i W_{u_{i}} \bar f_{v_{i}}\) with \(4\le u_{1}<v_{1}<\ldots<u_{k}<v_{k},\allowbreak \allowbreak k\ge 0\). The block's length intervals, the terminal intervals and $m$ may overlap. Within each block the written intervals are disjoint because the original residual table is strictly ordered. Equal-arm products \(E_t,\allowbreak \allowbreak E_r\) stay at their original vertices and satisfy \(0\le E_t,\allowbreak \allowbreak E_r\le 1\).

For each original connector length \(L\ge 1\) fix

\fitdisplay{
H_D=\begin{cases}
E_rP W_m T^L(E_tF),&m\le L,\\
E_rP \bar f_m A T^{L-1}(E_tF),&m>L.
\end{cases}
}{R36.1.1}

Every terminal \(h\) arm has its first edge toward \(t\), remaining edges toward the leaf; its \(f\) arm is entirely outward. \(W_m\) has its first edge outward and remaining edges inward toward that first internal point; \(\bar f_m\) is entirely inward. In the second row only the root-adjacent connector edge is A, from \(r\) toward \(t\); all remaining \(L-1\) connector edges are \(T\). The first row is all \(T\). Root \(P\) and equal pairs use their actual displayed directions.

Theorem A. Every original legal host satisfies

\fitdisplay{
\boxed{E H_D^8\le e^{-2d},\qquad d=\|K\one\|_2^2.}
}{R36.1.2}

The direction depends only on \(m,\allowbreak \allowbreak L\) and original arm tables. In particular \(L=1,\allowbreak \allowbreak 2\) always use the second row. No low/high host split selects directions. This theorem is about an ordinary eighth moment, not an arbitrary input or a same-direction reference-fourth or \(I_q\) assertion.

\paragraph{R36.2 Exact original finite tools}

All operators and inner products use the same \(\mu\). Put \(B=-K,\allowbreak \allowbreak u=B\one,\allowbreak \allowbreak  T=J+B,\allowbreak \allowbreak A=J-B=T^*,\allowbreak \allowbreak p=1-u,\allowbreak \allowbreak q=1+u,\allowbreak \allowbreak  f_n=A^n\one,\allowbreak \allowbreak \bar f_n=T^n\one,\allowbreak \allowbreak  h_n=Tf_{n-1},\allowbreak \allowbreak W_n=A\bar f_{n-1},\allowbreak \allowbreak \delta_i=\lVert B^i\one\rVert ^2,\allowbreak \allowbreak d=\delta_1,\allowbreak \allowbreak  \theta=\frac{13}{32},\allowbreak \allowbreak \kappa=\frac{19}{32},\allowbreak \allowbreak \tau=d-\delta_2\). C014 gives \(d\le \frac{1}{3},\allowbreak \allowbreak  \lVert B\rVert ^2\le \theta,\allowbreak \allowbreak  d^2\le \delta_2\le \theta\cdot d,\allowbreak \allowbreak  \delta_{i+1}\le \theta\cdot \delta_i\) and the original row \(\lvert Bv(x)\rvert \le \lVert v\rVert_2\). The \(f_n,\allowbreak \allowbreak \bar f_n\) means \(\mu_n\) lie in[0,1]; all four path types lie in[0,2]. The means $Eh_2=EW_2=1+d$ exceed one.

\(G_n(z)=\sum_{r=0}^{n-1}z^r,\allowbreak \allowbreak  G_{0}=0,\allowbreak \allowbreak  \beta_{1}=\beta_{2}=0\) and \(\beta_n=\sum_{i=1}^{{\lfloor (n-1)/2\rfloor}}(-1)^{i-1}\delta_i\). The actual finite identity \(1+G_n(B)u=\bar f_n+\sum_{r=3}^{n}\beta_r \bar f_{n-r}\) gives \(\tau\le \beta_n\le d\) for \(n\ge 3\). Finite geometry gives \(\lVert G_n\rVert \le 1\) for \(n=1\) or \(n\ge 3; G_{2}\) has square bound \(1+\theta=\frac{45}{32}\).

We use the exact signed exponential theorem, fully reconstructed in actual A013: for every original real \(V\),

\fitdisplay{
E e^{8BV}\le e^{8EBV+17\|V\|^2}.
}{R36.2.1}

It is the all-real fourth-exponent lemma applied to \(2V\); constant-vector parts and either sign of \(EBV\) are allowed.  C064, equation (9.21), gives \(\lVert T^n\rVert \le 1\) for every \(n\ge 3\). Exact C094's original \(T^2\) bound gives

\fitdisplay{
\|T^2\|\le(1+1024\delta_2/855)^2\le(\frac{2981}{2565})^2
<(\frac{7}{6})^2<7/5.
}{R36.2.2}

The rational comparison use \(s_{17886}<17955\) and \(245<252\). Thus for \(L\ge 3,\allowbreak \allowbreak  \lVert T^{L-1}\rVert \le \frac{7}{5}\), including the genuine \(L=3\) exception.

The strong ACTUAL unequal-pair envelopes, from C096, Section 2.2, are \(h_{4} f_b\le \exp(R_{4b}-69d/20),\allowbreak \allowbreak  h_a f_b\le \exp(R_{ab}-16d/5)\) for \(a\ge 5\), where \(R_{ab}=\sum_{n=a+1}^{b}K^n\one\). The \(a=4\) retained negative term \(-\tau f_{2}\) and \(f_{2}\ge p^2/2\) give a positive error at most \(\delta_2+\delta_2^2/(2\tau)\le 663d/1216\); subtract it from \(4d\) to get \(4201d/1216>69d/20\). For \(a\ge 5\) the actual finite mean \(\mu_{a-1}\le 1-83d/32\) and positive difference tail \(\le 4\delta_2/(1-\theta^2)\) give loss \(\frac{83}{16}-\frac{1664}{855}>\frac{16}{5}\). These are finite unequal-pair bounds on the actual arm messages.

\paragraph{R36.3 One terminal mask and its signed moments}

Write \(F\le \exp(R-\Gamma)\), where \(R\) is the one finite union of the terminal intervals. Always \(\Gamma\ge 16jd/5\); if \(a_{1}=4\) then \(\Gamma\ge 69d/20+16(j-1)d/5\), an extra \(d/4\) beyond \(16jd/5\). There is a finite 0-1 polynomial \(Q\) with \(R=-B^4 Q(-B)u\). \(At z=i \sqrt{x}\), write \(Q(-z)=U(x)-zV(x)\). Finite alternating subseries give

\fitdisplay{
\begin{gathered}
-\frac{x}{1-x^2}\le U,V\le\frac1{1-x^2},\quad\\
\lvert Q\rvert ^2\le\frac{1+x}{(1-x^2)^2}\le C_*:=32768/16245.
\end{gathered}
}{R36.3.1}

Only finite subseries occur; geometric tails bound them without changing an original path. With \(C_\beta=\frac{1024}{855}<\frac{6}{5}\) and \(\nu\) the original \(u\) spectral measure, \(R\) has coordinates \(-x^2U+z x^2V\). Let \(D_i=EB^iR,\allowbreak \allowbreak  \sigma=ER^2\). The exact identities and useful bounds are

\fitdisplay{
D_0=-\int x^2V,\quad D_1=\int x^2U,\quad D_2=\int x^3V,
}{R36.3.2}

\fitdisplay{
\begin{gathered}
\sigma\le C_*\theta^4d=28561d/519840<d/18,\\
\quad \lvert R\rvert ^2\le C_*\theta^3d<9d/64<\frac{1}{16},
\end{gathered}
}{R36.3.3}

\fitdisplay{
\begin{gathered}
\lvert D_0\rvert ,\lvert D_1\rvert <d/5,\qquad D_0\le C_\beta \delta_4<d/12,\\
\qquad D_0-D_1\le2C_\beta \delta_4<d/6.
\end{gathered}
}{R36.3.4}

Here \(C_\beta \theta^2=\frac{169}{855}<\frac{1}{5}\) and \(C_\beta \theta^3=\frac{2197}{27360}<\frac{1}{12}\). The last inequality in(R36.3.4) bounds \(U+V\) together. Also

\fitdisplay{
\begin{gathered}
D_0+D_1-D_2\le C_\beta(2\delta_3+\delta_4)<d/2,\qquad\\
\|TR\|^2=D_0^2+2D_0D_1+\|BR\|^2\\
\le(\frac{1}{25}+\frac{5}{216})d<(\frac{4}{15})^2d.
\end{gathered}
}{R36.3.5}

For the first use \(C_\beta<\frac{6}{5},\allowbreak \allowbreak \theta^2<\frac{1}{6},\allowbreak \allowbreak 2+\theta<\frac{5}{2}\). For the second retain the actual constant cross, then use \(d\le \frac{1}{3},\allowbreak \allowbreak  \lvert D_0\rvert ,\allowbreak \allowbreak \lvert D_1\rvert \le d/5\) and \(\theta \sigma<5d/216\). The final rational inequality is \(\frac{341}{5400}<\frac{16}{225}\). No constant vector has been replaced by \(u\).

If the terminal minimum is at least five, a sharper finite factorization is available: \(R=B^5 Q(-B)u=x^3V+z x^2U\). Then \(\sigma\le C_*\theta^5d<d/40\). The same bounds(R36.3.4) still hold, and

\fitdisplay{
-\langle G_2(B)u,R\rangle
=-\int x^3(U+V)\le2C_\beta \delta_5<d/15.
}{R36.3.6}

The final bound uses \(C_\beta<\frac{6}{5},\allowbreak \allowbreak \theta^4<\frac{1}{36}\). This improvement is needed only in the \(L=1\) proof when the first terminal arm is not four.

\paragraph{R36.4 True propagated perturbation, all original lengths}

Since \(\lvert R\rvert <\frac{1}{4}\), the scalar inequality \(e^R\le 1+R+(\frac{3}{5})R^2\) is valid; the complete \(\lvert R\rvert \le \frac{1}{2}\) factorial/integral certificate appears in actual A016, Section 4. The upper polynomial is positive. If \(S_L\) is \(T^L\) or \(A T^{L-1}\), positivity gives

\fitdisplay{
\begin{gathered}
0\le S_LF\le e^{-\Gamma}\{S_L1+S_LR+\omega_L\},\\
\quad\omega_1=(\frac{6}{5})\sigma<d/15,\\
\quad\omega_L=(\frac{12}{5})\sigma<2d/15\ (L\ge2).
\end{gathered}
}{R36.4.1}

Indeed the last positive \(A/T\) row is at most twice the preceding mean; for \(L\ge 2\) that mean is \(E[f_{L-1}R^2]\le 2\sigma\). \(At L=1\) it is \(\sigma\). The braced expression is a nonnegative upper for the complete sum; its summands can have either sign.

Set \(V_R=T^{L-1}R\) and \(a_R=EV_R\). Then \(S_LR=a_R+B V_R\) for \(S_L=T^L\), and \(a_R-B V_R\) for \(S_L=A T^{L-1}\). For \(L=1,\allowbreak \allowbreak  EA R=D_0-D_1<d/6\). For \(L=2,\allowbreak \allowbreak  EA TR=D_0+D_1-D_2<d/2\) and \(\lVert V_R\rVert <4\sqrt{d}/15\) by(R36.3.5).

For all \(L\ge 3,\allowbreak \allowbreak  (\) R36.2.2) and(R36.3.3) give \(\lVert V_R\rVert <7\sqrt{d}/20\). There is also a uniform true mean bound. The exact finite Duhamel expansion writes, for every \(k\ge 1\),

\fitdisplay{
f_k=\mu_{k-1}1+\sum_{j=1}^{k-1}\mu_{k-1-j}(-B)^j\one+(-B)^k\one.
}{R36.L248}

Every coefficient is in[0,1]. Consequently \(\lVert f_k-\mu_{k-1}1\rVert \le \sqrt{d}\sum_{j=0}^{k-1}\theta^{j/2}<3\sqrt{d}\). Hence \(a_R=\langle f_{L-1},\allowbreak \allowbreak R\rangle \le d/12+3\sqrt{d \sigma}<5d/6\). Skewness gives \(\lvert EBV_R\rvert \le \sqrt{d}\lVert V_R\rVert \). Both possible last-edge signs therefore satisfy

\fitdisplay{
E S_LR<71d/60\qquad(L\ge3).
}{R36.4.2}

\paragraph{R36.5 Actual bare sums and one eighth-moment payment formula}

\(At L=1\) the bare sum is \(\bar f_m+p,\allowbreak \allowbreak  m\ge 3\). The exact finite convolution gives

\fitdisplay{
\begin{gathered}
\bar f_m+p\le 2-\beta_m+B V_b, V_b=G_{m-1}(B)u,\\
EBV_b=-\beta_{m+1}, \beta_m+\beta_{m+1}\ge 2\tau\ge 19d/16.
\end{gathered}
}{R36.L265}

\(At L=2\) it is \(\bar f_m+W_{2}\). Here

\fitdisplay{
\begin{gathered}
\bar f_m+W_{2}\le 2-\beta_m+B V_b, V_b=B G_{m-2}(B)u,\\
EBV_b\le \delta_2, \lVert V_b\rVert \le 4\sqrt{d}/5,\\
-8\beta_m+8EBV_b\le -3d/2.
\end{gathered}
}{R36.5.1}

The \(G_{2}\) exception is paid because \(\theta(1+\theta)=\frac{585}{1024}<\frac{16}{25}\). For \(L\ge 3\) the structure-only direction in(R36.1.1) makes the sum \(W_h+\bar f_H,\allowbreak \allowbreak  h=\min(m,\allowbreak \allowbreak L)\ge 3,\allowbreak \allowbreak  H=\max(m,\allowbreak \allowbreak L)\). Its exact finite pair estimate, reconstructed in R2-C030/A014, is

\fitdisplay{
\begin{gathered}
W_h+\bar f_H\le2-(\frac{51}{32})d+B V_b,\quad\\
EBV_b\le\theta d,\quad\|V_b\|\le\tfrac12\sqrt d.
\end{gathered}
}{R36.5.2}

For \(h<H,\allowbreak \allowbreak  V_b=B^{h-1}G_{H-h}(B)u; h=H\) has \(V_b=0\) and \(W_h+\bar f_h=2\mu_{h-1}\le 2-2d\). The finite \(G_{2}\) exception gives \(\lVert V_b\rVert ^2\le \theta^2(1+\theta)d<d/4\); the first positive finite mean is \(\delta_2\). Thus (R36.5.2) includes every finite \(m,\allowbreak \allowbreak L\) and equality \(m=L\).

Combine the appropriate bare sum, the same true perturbation, and(R36.4.1). If its bare loss is \(\lambda\cdot d (\) or \(\beta_m\)), actual nonnegative AM--GM and (R36.2.1) give the unpowered-root-pair estimate

\fitdisplay{
\begin{gathered}
\log E H^8\le-8\Gamma-8\lambda d+8EBV_b\\
+8\omega_L+8E S_LR+17\|V_b\mathbin{\pm}V_R\|^2.
\end{gathered}
}{R36.5.3}

The sign is minus in the A row and plus in the \(T\) row. The right side comes from a pointwise positive upper for the actual sum.

\paragraph{R36.6 All \(L\) without root post-pairs}

\paragraph{R36.6.1 \(L=1\), terminal minimum exactly four}

For \(m=3,\allowbreak \allowbreak  V_b=G_2 u\). From \(R=-B^4 Q(-B)u\),

\fitdisplay{
-\langle G_2 u,R\rangle =\int [x^2U-x^3V]\le C_\beta(\delta_3+\delta_5)<7d/30.
}{R36.L308}

Therefore

\fitdisplay{
\|V_b-R\|^2\le(\frac{45}{32}+\frac{7}{15}+\frac{1}{18})d=(\frac{2777}{1440})d.
}{R36.6.1}

For \(m\ge 4,\allowbreak \allowbreak  G_{m-1}\) contracts, so the weaker bound \(\lVert V_b-R\rVert ^2<(\frac{5}{4})^2d\) is even smaller than(R36.6.1). Using only \(\Gamma\ge 16d/5\) first, (R36.5.3) gives a loss coefficient at least

\fitdisplay{
\frac{128}{5}+\frac{19}{2}-17(\frac{2777}{1440})-\frac{8}{15}-\frac{4}{3}=647/1440.
}{R36.L320}

\(At\) actual \(a_{1}=4\) the retained \(\frac{69}{20}\) first-pair loss adds 2 to this coefficient. Thus the full \(L=1\) loss exceeds 2, for all \(m\ge 3\) and all \(j\).

\paragraph{R36.6.2 \(L=1\), terminal minimum at least five}

For \(m=3\) use(R36.3.6), yielding \(\lVert V_b-R\rVert ^2\le (\frac{45}{32}+\frac{2}{15}+\frac{1}{40})d=(\frac{751}{480})d\). For \(m\ge 4\) the bound \(25d/16=750d/480\) suffices. Even keeping the weaker \(\omega_1<d/15\) and mean \(E S_1 R<d/6\), the coefficient is at least

\fitdisplay{
\frac{128}{5}+\frac{19}{2}-17(\frac{751}{480})-\frac{8}{15}-\frac{4}{3}=\frac{637}{96}>2.
}{R36.6.2}

The two structural cases give the bound for \(a_{1}\ge 4\).

\paragraph{R36.6.3 \(L=2\)}

By(R36.3.5) and(R36.5.1 \(),\allowbreak \allowbreak  \lVert V_b-V_R\rVert <16\sqrt{d}/15\). Use \(\Gamma\ge 16d/5,\allowbreak \allowbreak  \omega_2<2d/15\) and \(E S_2 R<d/2\). Formula(R36.5.3) gives

\fitdisplay{
\frac{128}{5}+\frac{3}{2}-17(\frac{16}{15})^2-\frac{16}{15}-4=\frac{1211}{450}>2.
}{R36.6.3}

This uses the genuine constant-vector square \(\lVert TR\rVert ^2\) in(R36.3.5).

\paragraph{R36.6.4 Every \(L\ge 3\)}

Equations(R36.4.2),(R36.5.2) give \(\lVert V_b\pm V_R\rVert <17\sqrt{d}/20\). The norm cost \(17(\frac{17}{20})^2<13\), and the combined \(\delta/\) mean fee is \(\frac{16}{15}+\frac{142}{15}=\frac{158}{15}<11\). The loss before those fees is \(\frac{128}{5}+8(\frac{51}{32}-\theta)=\frac{351}{10}\). Therefore the remaining coefficient exceeds \(\frac{351}{10}-13-11=\frac{111}{10}>2\). Every finite \(L\) is covered by the actual power bound and finite path identity, including \(L=3\) separately in(R36.2.2).

\paragraph{R36.7 All finite negative root post-pairs, including overlapping masks}

For \(k\ge 1\), the actual \(P=\prod W_{u_{i}} \bar f_{v_{i}}\) has the reversed same pair envelope \(P\le \exp(S-\Gamma_r),\allowbreak \allowbreak  \Gamma_r\ge 16kd/5\). Write \(S=B V_p=B^4  Q_p(B)u\), where \(Q_p\) is one finite 0-1 mask. The same finite-coordinate proof as(R36.3.1) gives

\fitdisplay{
ES<d/12,\qquad\|V_p\|<3\sqrt d/8.
}{R36.7.1}

The ordinary exponent for \(H\) now contains the SINGLE actual vector \(V_b\pm V_R+V_p\), with additional loss \(8\Gamma_r\) and mean \(8ES\). No assertion that terminal and root masks are disjoint has been used. Their cross term is paid by the vector triangle bound below.

\(L=1: (\) R36.6.1) gives \(\lVert V_b-R\rVert <7\sqrt{d}/5\), so the combined norm is less than \(71\sqrt{d}/40\). Its 17-times square is \(<54\). The \(\delta\), perturbation mean and post mean fees sum to \(\frac{8}{15}+\frac{4}{3}+\frac{2}{3}=\frac{38}{15}<3\). Even at \(j=k=1\), the loss \(8\Gamma+8\Gamma_r\) plus the bare signed loss is at least \(607d/10\). The final coefficient therefore exceeds \(\frac{607}{10}-54-3=\frac{37}{10}>2\).

\(L=2\): the combined norm is at most \(173\sqrt{d}/120\), whose 17-times square is \(<36\). The other fees \(\frac{16}{15}+4+\frac{2}{3}=\frac{86}{15}\) are \(<6\). The starting loss is \(\frac{256}{5}+\frac{3}{2}=\frac{527}{10}\). The remaining coefficient exceeds \(\frac{107}{10}>2\).

\(L\ge 3\): the combined norm is at most \(49\sqrt{d}/40\), whose 17-times square is \(<26\). Other fees \(\frac{16}{15}+\frac{142}{15}+\frac{2}{3}=\frac{56}{5}<12\). Starting loss \(\frac{607}{10}\) leaves a coefficient greater than \(\frac{227}{10}>2\).

Extra terminal or root pairs only increase \(\Gamma\) or \(\Gamma_r\); the finite mask estimates already hold for every finite pair count. Thus all \(k\ge 1\) and \(j\ge 1\) are paid with the SAME bound \(e^{-2d}\). Restore \(E_t,\allowbreak \allowbreak E_r\) by the actual positive original operators and multiplications to prove(R36.1.2).

